% ============================================================
% Part 1 — Basic Worked Problems
% Topics: language of polynomials, simple factor/remainder theorem,
%         straightforward Vieta's, basic graph sketching
% ============================================================

\begin{problem}[Identically Equal Polynomials]
Find $a$, $b$ and $c$ given that
\[
a(x-2)^2+b(x-2)+c=x^2+x+1
\quad\text{for all }x.
\]
\end{problem}

\begin{hintbox}
Expand the left-hand side and compare coefficients.
\end{hintbox}

\begin{solution}
\[
a(x-2)^2+b(x-2)+c=ax^2+(-4a+b)x+(4a-2b+c).
\]
Comparing with $x^2+x+1$ gives
\[
a=1,\qquad -4a+b=1,\qquad 4a-2b+c=1.
\]
Hence
\[
\boxed{a=1,\quad b=5,\quad c=7.}
\]
\end{solution}

\begin{takeaways}
\begin{itemize}[leftmargin=*]
\item Polynomial identities are solved by coefficient comparison.
\item Expanding around $(x-2)$ is still just an ordinary polynomial identity.
\end{itemize}
\end{takeaways}

\begin{problem}[Odd Polynomial Coefficients]
Suppose
\[
P(x)=ax^5+bx^4+cx^3+dx^2+ex+f
\]
is an odd function, so $P(-x)=-P(x)$ for all $x$. Show that $b=d=f=0$.
\end{problem}

\begin{hintbox}
Write down $P(-x)$ and compare it with $-P(x)$ term by term.
\end{hintbox}

\begin{solution}
\[
P(-x)=-ax^5+bx^4-cx^3+dx^2-ex+f,
\]
while
\[
-P(x)=-ax^5-bx^4-cx^3-dx^2-ex-f.
\]
Since these are identical,
\[
b=-b,\qquad d=-d,\qquad f=-f.
\]
So
\[
\boxed{b=d=f=0.}
\]
\end{solution}

\begin{takeaways}
\begin{itemize}[leftmargin=*]
\item Odd polynomials contain only odd powers of $x$.
\item The constant term of an odd polynomial must be $0$.
\end{itemize}
\end{takeaways}

\begin{problem}[Central Coefficient]
What is the coefficient of $x^n$ in
\[
G(x)=\left(1+x+x^2+\cdots+x^n\right)^2?
\]
\end{problem}

\begin{hintbox}
Count the pairs $(i,j)$ with $0\le i,j\le n$ and $i+j=n$.
\end{hintbox}

\begin{solution}
Write
\[
G(x)=\left(\sum_{i=0}^{n}x^i\right)\left(\sum_{j=0}^{n}x^j\right).
\]
The term $x^n$ appears whenever $x^i\cdot x^j=x^n$, that is, whenever $i+j=n$.
The possible pairs are
\[
(0,n),(1,n-1),\ldots,(n,0),
\]
which gives $n+1$ choices. Therefore the coefficient is
\[
\boxed{n+1.}
\]
\end{solution}

\begin{takeaways}
\begin{itemize}[leftmargin=*]
\item Coefficient questions in products often reduce to counting exponent pairs.
\item The coefficient of $x^r$ is the number of ways to split $r$ as a sum of exponents.
\end{itemize}
\end{takeaways}

\begin{problem}[Easy Zeroes First]
Solve
\[
x^3-x^2-4x+4=0
\]
by first factoring the left-hand side.
\end{problem}

\begin{hintbox}
Factor by grouping.
\end{hintbox}

\begin{solution}
\[
x^3-x^2-4x+4=x^2(x-1)-4(x-1)=(x-1)(x^2-4).
\]
So
\[
(x-1)(x-2)(x+2)=0.
\]
Hence
\[
\boxed{x=1,\ 2,\ -2.}
\]
\end{solution}

\begin{takeaways}
\begin{itemize}[leftmargin=*]
\item Before using heavier tools, try grouping or simple integer roots.
\item Cubics designed for students often hide an obvious factor.
\end{itemize}
\end{takeaways}

\begin{problem}[Factors of $x^n+1$]
Is either $x+1$ or $x-1$ a factor of $x^n+1$, where $n$ is a positive integer?
\end{problem}

\begin{hintbox}
Use the factor theorem at $x=1$ and $x=-1$.
\end{hintbox}

\begin{solution}
Let $P(x)=x^n+1$.

For $x-1$ to be a factor we need $P(1)=0$, but
\[
P(1)=2\ne 0.
\]
So $x-1$ is never a factor.

For $x+1$ to be a factor we need $P(-1)=0$:
\[
P(-1)=(-1)^n+1.
\]
If $n$ is odd, this is $0$. If $n$ is even, this is $2$.

Therefore
\[
\boxed{x+1\text{ is a factor iff }n\text{ is odd, and }x-1\text{ is never a factor}.}
\]
\end{solution}

\begin{takeaways}
\begin{itemize}[leftmargin=*]
\item The factor theorem turns factor questions into substitutions.
\item Powers of $-1$ almost always split into odd and even cases.
\end{itemize}
\end{takeaways}

\begin{problem}[Simple Partial Fractions]
Given
\[
\frac{1}{x^3-1}=\frac{A}{x-1}+\frac{Bx+C}{x^2+x+1},
\]
find $A$, $B$ and $C$.
\end{problem}

\begin{hintbox}
Use $x^3-1=(x-1)(x^2+x+1)$, then clear denominators.
\end{hintbox}

\begin{solution}
Since
\[
x^3-1=(x-1)(x^2+x+1),
\]
we have
\[
1=A(x^2+x+1)+(Bx+C)(x-1).
\]
Expand:
\[
1=(A+B)x^2+(A-B+C)x+(A-C).
\]
Comparing coefficients gives
\[
A+B=0,\qquad A-B+C=0,\qquad A-C=1.
\]
Solving,
\[
\boxed{A=\frac13,\qquad B=-\frac13,\qquad C=-\frac23.}
\]
\end{solution}

\begin{takeaways}
\begin{itemize}[leftmargin=*]
\item Clear denominators before comparing coefficients.
\item Partial fractions are another good place to use polynomial identities.
\end{itemize}
\end{takeaways}
